\section{模型假设}

为建立清晰且可求解的模型，本文作出以下假设：

\begin{enumerate}
  \item \textbf{纵向线性增长假设}：同一孕妇的 Y 染色体浓度在检测窗口内随孕周近似线性上升。数据中 260 名多次采血孕妇中 226 名（86.9\%）末期浓度高于初期，支持该假设。
  \item \textbf{BMI 稀释效应假设}：孕妇 BMI 越高，母体血液循环量越大，对胎儿游离 DNA 的稀释作用越强，导致 Y 染色体浓度越低。该假设符合 NIPT 的生理机制，并在数据中得到验证。
  \item \textbf{共享斜率假设}：不同孕妇的 Y 染色体浓度随孕周的增长率相近，个体差异主要体现在浓度水平（截距）上。据此可用统一的固定效应斜率估计个体达标时间。
  \item \textbf{正态误差假设}：在控制孕周与 BMI 后，Y 染色体浓度的随机误差近似服从正态分布，这是混合效应模型与达标概率计算的基础。
  \item \textbf{达标阈值假设}：以题面给定的 $4\%$ 作为 Y 染色体浓度达标阈值，$Y\ge 4\%$ 时 NIPT 结果基本准确，否则难以保证。
  \item \textbf{风险分段假设}：晚检风险按题面分段——早期（$\le 12$ 周）低风险、中期（13--27 周）高风险、晚期（$\ge 28$ 周）极高风险。
  \item \textbf{标签一致性假设}：女胎异常判定以附件 AB 列（染色体非整倍体，空白即无异常）为判定结果标签。
  \item \textbf{检测误差独立假设}：检测误差为独立的随机噪声，其标准差记为 $\sigma_e$，不随孕周与 BMI 系统变化。
\end{enumerate}

上述假设的必要性与影响范围将在相应模型处进一步说明。
